% !TeX program = lualatex
% Conversão do documento Word 08_CV_PCV-PTV_2025.docx para LaTeX.
% Compilar preferencialmente com LuaLaTeX ou XeLaTeX.
\documentclass[12pt,letterpaper]{article}

\usepackage{fontspec}
\IfFontExistsTF{Trebuchet MS}{\setmainfont{Trebuchet MS}}{\setmainfont{Noto Sans}}
\IfFontExistsTF{Microsoft Sans Serif}{\newfontfamily\msfont{Microsoft Sans Serif}}{\newfontfamily\msfont{Noto Sans}}
\newfontfamily\trebfont{Noto Sans}
\IfFontExistsTF{Trebuchet MS}{\renewfontfamily\trebfont{Trebuchet MS}}{\renewfontfamily\trebfont{Noto Sans}}

\usepackage[brazil]{babel}
\usepackage{amsmath,amssymb}
\usepackage{graphicx}
\usepackage{float}
\usepackage{geometry}
\usepackage{setspace}
\usepackage{titlesec}
\usepackage{enumitem}
\usepackage{changepage}
\usepackage{ragged2e}
\usepackage{microtype}
\usepackage{caption}
\usepackage{booktabs}
\usepackage{array}

% Configuração de página observada no Word: Letter, margens aprox. 2,54 cm laterais,
% 3,00 cm superior e 1,76 cm inferior.
\geometry{letterpaper,left=2.54cm,right=2.54cm,top=3.00cm,bottom=1.76cm}

\setlength{\parindent}{0pt}
\setlength{\parskip}{6pt}
\setstretch{1.05}
\justifying

% Estilos de títulos aproximados aos estilos do Word.
\titleformat{\section}
  {\trebfont\fontsize{18}{22}\selectfont\bfseries\justifying}{}{0pt}{}
\titlespacing*{\section}{0pt}{3.8pt}{8pt}

\titleformat{\subsection}
  {\msfont\fontsize{17}{21}\selectfont\bfseries\justifying}{}{0pt}{}
\titlespacing*{\subsection}{0pt}{4pt}{6pt}

\titleformat{\subsubsection}
  {\trebfont\fontsize{16}{20}\selectfont\bfseries\justifying}{}{0pt}{}
\titlespacing*{\subsubsection}{0pt}{4pt}{5pt}

\titleformat{\paragraph}
  {\msfont\fontsize{14}{18}\selectfont\bfseries}{}{0pt}{}
\titlespacing*{\paragraph}{0pt}{5pt}{4pt}

\setlist[itemize]{label=\textbullet,leftmargin=1.25cm,itemsep=2pt,topsep=4pt}
\setlist[enumerate]{leftmargin=1.25cm,itemsep=2pt,topsep=4pt}

\newenvironment{corpo}{\begin{adjustwidth}{0.46cm}{0.45cm}\justifying}{\end{adjustwidth}}
\newcommand{\m}[1]{\ensuremath{#1}}

\begin{document}

\section*{Relacionando as cotas de A (PCV) e B (PTV) a partir do PIV}

\begin{corpo}
Seja uma parábola de concordância vertical de comprimento $L$, com tangentes de inclinação $i_1$ (em A) e $i_2$ (em B), sendo o PIV (I) o ponto médio da curva ($x_I = L/2$). Sabendo a cota $Z_I$ do PIV, podemos determinar as cotas dos pontos $A$ (PCV) e $B$ (PTV) (Figura):

\begin{figure}[H]
\centering
\includegraphics[width=0.92\linewidth]{figuras/figura_parabola_piv_pcv_ptv.png}
\end{figure}
\end{corpo}

\subsubsection*{Equações gerais}

\begin{corpo}
A equação da parábola (com origem em $A$) é:
\[
 y(x)=i_1x-\frac{(i_1-i_2)x^2}{2L}
\]

A cota absoluta ao longo da parábola:
\[
 Z(x)=Z_A+y(x)
\]

Cota do PIV ($x_I=L/2$):
\[
\begin{aligned}
Z_I
&=Z_A+y\left(\frac{L}{2}\right) \\
&=Z_A+i_1\frac{L}{2}-\frac{(i_1-i_2)}{2L}\left(\frac{L}{2}\right)^2 \\
&=Z_A+i_1\frac{L}{2}-\frac{(i_1-i_2)L}{8}.
\end{aligned}
\]

\paragraph*{Isolando $Z_A$:}
\[
 Z_A=Z_I-i_1\frac{L}{2}+\frac{(i_1-i_2)L}{8}
\]

Cota de $B$ (PTV, $x=L$):
\[
\begin{aligned}
y(L)&=i_1L-\frac{(i_1-i_2)L^2}{2L} \\
&=i_1L-\frac{(i_1-i_2)L}{2} \\
&=\frac{(i_1+i_2)L}{2}.
\end{aligned}
\]

Portanto,
\[
 Z_B=Z_A+\frac{(i_1+i_2)L}{2}
\]

Ou, substituindo $Z_A$:
\[
 Z_B=Z_I-i_1\frac{L}{2}+\frac{(i_1-i_2)L}{8}+\frac{(i_1+i_2)L}{2}
\]

\begin{figure}[H]
\centering
\includegraphics[width=0.92\linewidth]{figuras/figura_perfil_longitudinal_convexa.jpeg}
\end{figure}
\end{corpo}

\section*{Relação com a flecha vertical $e$ e da cota do PIV}

\begin{corpo}
Vamos repetir a demonstração anterior. Para isso, considere novamente uma curva de concordância vertical do tipo parábola simples, de comprimento $L$, com rampas tangentes de inclinação $i_1$ (em A) e $i_2$ (em B), e $I$ (PIV) sendo o ponto médio da curva ($x_I=L/2$). A flecha vertical $e$ pode ser calculada por:
\[
 e=\frac{L}{8}(i_1-i_2)
\]

A equação da parábola (com origem em $A$) é:
\[
 y(x)=i_1x-\frac{(i_1-i_2)x^2}{2L}
\]

A cota absoluta ao longo da parábola é:
\[
 Z(x)=Z_A+y(x)
\]
\end{corpo}

\subsubsection*{Cota do ponto $I$ (PIV)}

\begin{corpo}
Para $x_I=L/2$, temos:
\[
\begin{aligned}
y\left(\frac{L}{2}\right)
&=i_1\frac{L}{2}-\frac{(i_1-i_2)}{2L}\left(\frac{L}{2}\right)^2 \\
&=i_1\frac{L}{2}-\frac{(i_1-i_2)L}{8} \\
&=i_1\frac{L}{2}-e.
\end{aligned}
\]

Logo,
\[
 Z_I=Z_A+i_1\frac{L}{2}-e
\]

Isolando $Z_A$ em função de $Z_I$ e $e$:
\[
 Z_A=Z_I-i_1\frac{L}{2}+e
\]
\end{corpo}

\subsubsection*{Cota do ponto $B$ (PTV)}

\begin{corpo}
Para $x=L$:
\[
\begin{aligned}
y(L)&=i_1L-\frac{(i_1-i_2)L^2}{2L} \\
&=i_1L-\frac{(i_1-i_2)L}{2} \\
&=\frac{(i_1+i_2)L}{2}.
\end{aligned}
\]

Portanto,
\[
 Z_B=Z_A+\frac{(i_1+i_2)L}{2}
\]

Substituindo $Z_A$:
\[
 Z_B=Z_I-i_1\frac{L}{2}+e+\frac{(i_1+i_2)L}{2}
\]
\end{corpo}

\subsubsection*{Resumo das Equações}

\begin{corpo}
\[
 Z_A=Z_I-i_1\frac{L}{2}+e
\]
\[
 Z_B=Z_I-i_1\frac{L}{2}+e+\frac{(i_1+i_2)L}{2}
\]

onde
\[
 e=\frac{L}{8}(i_1-i_2).
\]

As duas formas de calcular as cotas de A e B nos levam ao mesmo resultado, sendo apenas explicitado em termos de $e$ a segunda opção.
\end{corpo}

\subsubsection*{Exemplo Numérico}

\begin{corpo}
Considere:
\begin{itemize}
    \item $i_1=+2{,}5\%=0{,}025$;
    \item $i_2=-1{,}0\%=-0{,}010$;
    \item $L=120\,\text{m}$;
    \item $Z_I=201{,}20\,\text{m}$.
\end{itemize}

Cálculo de $Z_A$:
\[
\begin{aligned}
Z_A
&=201{,}20-0{,}025\cdot 60+\frac{(0{,}025-(-0{,}010))\cdot 120}{8} \\
&=201{,}20-1{,}50+\frac{0{,}035\cdot 120}{8} \\
&=201{,}20-1{,}50+\frac{4{,}20}{8} \\
&=201{,}20-1{,}50+0{,}525 \\
&=200{,}225\,\text{m}.
\end{aligned}
\]

Cálculo de $Z_B$:
\[
\begin{aligned}
Z_B
&=201{,}20-0{,}025\cdot 60+\frac{(0{,}025-(-0{,}010))\cdot 120}{8}
  +\frac{(0{,}025+(-0{,}010))\cdot 120}{2} \\
&=201{,}20-1{,}50+0{,}525+\frac{0{,}015\cdot 120}{2} \\
&=201{,}20-1{,}50+0{,}525+0{,}90 \\
&=201{,}20-1{,}50+1{,}425 \\
&=201{,}125\,\text{m}.
\end{aligned}
\]

Resumo:
\begin{itemize}
    \item $Z_A=200{,}225\,\text{m}$;
    \item $Z_I=201{,}20\,\text{m}$;
    \item $Z_B=201{,}125\,\text{m}$.
\end{itemize}
\end{corpo}

\section*{Simetria das Cotas de A (PCV) e B (PTV) para $|i_1|=|i_2|$}

\begin{corpo}
Suponha que as inclinações das rampas sejam simétricas em módulo, ou seja, $i_1=-i_2$, com $i_1>0$ e $i_2<0$. A equação geral para as cotas dos pontos extremos, obtida anteriormente, é:
\[
 Z_A=Z_I-\frac{i_1L}{2}+\frac{(i_1-i_2)L}{8}
\]
\[
 Z_B=Z_A+\frac{(i_1+i_2)L}{2}
\]

Substituindo $i_2=-i_1$:
\[
 i_1-i_2=i_1-(-i_1)=2i_1
\]
\[
 i_1+i_2=i_1+(-i_1)=0
\]

Assim,
\[
\begin{aligned}
Z_A
&=Z_I-\frac{i_1L}{2}+\frac{2i_1L}{8} \\
&=Z_I-\frac{i_1L}{2}+\frac{i_1L}{4} \\
&=Z_I-\frac{2i_1L}{4}+\frac{i_1L}{4} \\
&=Z_I-\frac{i_1L}{4}.
\end{aligned}
\]

\[
 Z_B=Z_A+0=Z_A
\]
\end{corpo}

\subsubsection*{Conclusão}

\begin{corpo}
Quando $i_1=-i_2$, as cotas de A (PCV) e B (PTV) são sempre iguais:
\[
 Z_B=Z_A=Z_I-\frac{i_1L}{4}
\]
\end{corpo}

\subsubsection*{Exemplo Numérico}

\begin{corpo}
Sejam:
\begin{itemize}
    \item $i_1=+2{,}5\%=0{,}025$;
    \item $i_2=-2{,}5\%=-0{,}025$;
    \item $L=100\,\text{m}$;
    \item $Z_I=250{,}00\,\text{m}$.
\end{itemize}

\[
\begin{aligned}
Z_A
&=250{,}00-\frac{0{,}025\times 100}{4} \\
&=250{,}00-\frac{2{,}5}{4} \\
&=250{,}00-0{,}625 \\
&=249{,}375\,\text{m}.
\end{aligned}
\]

\[
 Z_B=Z_A=249{,}375\,\text{m}
\]

Portanto, neste caso simétrico, a cota inicial $Z_A$ e a cota final $Z_B$ são exatamente iguais.
\end{corpo}

\section*{Vértice da parábola de concordância vertical}

\begin{corpo}
Considere a parábola de concordância vertical, definida por:
\[
 y(x)=i_1x-\frac{(i_1-i_2)x^2}{2L}
\]

onde:
\begin{itemize}
    \item $i_1$ = inclinação inicial (em $x=0$);
    \item $i_2$ = inclinação final (em $x=L$);
    \item $L$ = comprimento da curva.
\end{itemize}

O ponto de máximo ou mínimo da parábola ocorre no vértice $(x_0,y_0)$.
\end{corpo}

\subsubsection*{Cálculo do vértice}

\begin{corpo}
\textbf{1. Valor de $x_0$ (posição do vértice):}

O vértice é obtido anulando a derivada de $y(x)$:
\[
\frac{dy}{dx}=i_1-\frac{(i_1-i_2)}{L}x=0
\]
\[
\frac{(i_1-i_2)}{L}x_0=i_1
\]
\[
 x_0=\frac{i_1L}{i_1-i_2}
\]

\textbf{2. Valor de $y_0$ (valor da parábola no vértice):}
\[
\begin{aligned}
y_0&=y(x_0)=i_1x_0-\frac{(i_1-i_2)}{2L}x_0^2 \\
&=i_1\left(\frac{i_1L}{i_1-i_2}\right)
-\frac{(i_1-i_2)}{2L}\left(\frac{i_1L}{i_1-i_2}\right)^2 \\
&=\frac{i_1^2L}{i_1-i_2}-\frac{i_1^2L}{2(i_1-i_2)} \\
&=\frac{i_1^2L}{2(i_1-i_2)}.
\end{aligned}
\]

Resumo final:
\[
 x_0=\frac{i_1L}{i_1-i_2}
\]
\[
 y_0=\frac{i_1^2L}{2(i_1-i_2)}
\]
\end{corpo}

\section*{Vértice da parábola: posição e cota absoluta}

\begin{corpo}
Considere a parábola de concordância vertical para alinhamentos em projetos de vias, com equação relativa:
\[
 y(x)=i_1x-\frac{(i_1-i_2)x^2}{2L}
\]

onde:
\begin{itemize}
    \item $i_1$: inclinação da rampa tangente no ponto $A$ (PCV, $x=0$);
    \item $i_2$: inclinação da rampa tangente no ponto $B$ (PTV, $x=L$);
    \item $L$: comprimento da curva de concordância.
\end{itemize}

A cota absoluta em qualquer ponto da parábola é:
\[
 Z(x)=Z_A+y(x)
\]

onde $Z_A$ é a cota inicial no ponto $A$ (PCV).
\end{corpo}

\subsubsection*{Coordenadas do vértice da parábola (máximo ou mínimo)}

\begin{corpo}
O vértice, $(x_0,y_0)$, ocorre onde a derivada de $y(x)$ é nula:
\[
\frac{dy}{dx}=i_1-\frac{(i_1-i_2)}{L}x
\]
\[
 x_0=\frac{i_1L}{i_1-i_2}
\]

O valor de $y$ nesse ponto é:
\[
\begin{aligned}
y_0&=y(x_0)=i_1x_0-\frac{(i_1-i_2)}{2L}x_0^2 \\
&=i_1\left(\frac{i_1L}{i_1-i_2}\right)
-\frac{(i_1-i_2)}{2L}\left(\frac{i_1L}{i_1-i_2}\right)^2 \\
&=\frac{i_1^2L}{i_1-i_2}-\frac{i_1^2L}{2(i_1-i_2)} \\
&=\frac{i_1^2L}{2(i_1-i_2)}.
\end{aligned}
\]

A cota absoluta do vértice é:
\[
 Z_0=Z_A+y_0=Z_A+\frac{i_1^2L}{2(i_1-i_2)}
\]
\end{corpo}

\subsubsection*{Relação com PIV (I) e extremos da curva}

\begin{corpo}
\begin{itemize}
    \item O PIV (I) — ponto de interseção das tangentes das rampas — costuma ser localizado em $x=L/2$ (curva simétrica), mas o vértice da parábola, em geral, não coincide com o PIV.
    \item Os pontos extremos são:
\end{itemize}

Em $A$ (PCV, $x=0$):
\[
 Z_A
\]

Em $B$ (PTV, $x=L$):
\[
 Z_B=Z_A+\frac{(i_1+i_2)L}{2}
\]

\begin{itemize}
    \item O vértice $(x_0,Z_0)$ pode estar entre $A$ e $B$ ($0<x_0<L$) ou fora desse intervalo, dependendo das inclinações.
\end{itemize}
\end{corpo}

\subsubsection*{Resumo das fórmulas principais}

\begin{corpo}
\[
 x_0=\frac{i_1L}{i_1-i_2}
\]
\[
 y_0=\frac{i_1^2L}{2(i_1-i_2)}
\]
\[
 Z_0=Z_A+\frac{i_1^2L}{2(i_1-i_2)}
\]
\end{corpo}

\subsubsection*{Observações}

\begin{corpo}
\begin{itemize}
    \item Se $i_1>i_2$, a parábola é convexa e o vértice é um máximo.
    \item Se $i_1<i_2$, a parábola é côncava e o vértice é um mínimo.
    \item Se $x_0$ está fora do intervalo $[0,L]$, o máximo/mínimo ocorre em um dos extremos.
    \item O PIV (I), na forma simétrica usual, ocorre em $x=L/2$, e sua cota é $Z_I=Z_A+y(L/2)$.
\end{itemize}

Outra forma de analisar:
\begin{itemize}
    \item Se $0<x_0<L$, o vértice está entre $A$ e $B$.
    \item Se $i_1>i_2$, a parábola é convexa e $y_0$ é máximo.
    \item Se $i_1<i_2$, a parábola é côncava e $y_0$ é mínimo.
    \item Se $i_1=i_2$, a curva é uma reta.
\end{itemize}
\end{corpo}

\section*{Equações das retas tangentes em A (PCV) e B (PTV)}

\begin{corpo}
\begin{itemize}
    \item Reta tangente em $A$ (PCV):
\end{itemize}
\[
 Z_{\mathrm{tang},A}(x)=Z_A+i_1x
\]

onde $x$ é a abscissa a partir de $A$ e $Z_A$ a cota de $A$.

\begin{itemize}
    \item Reta tangente em $B$ (PTV):
\end{itemize}
\[
 Z_{\mathrm{tang},B}(x)=Z_B+i_2(x-L)
\]

onde $x$ é a abscissa ao longo do eixo do projeto, $L$ o comprimento da curva e $Z_B$ a cota de $B$.
\end{corpo}

\section*{Demonstração do ponto de interseção das tangentes}

\begin{corpo}
As equações das tangentes são:
\[
 Z_{\mathrm{tang},A}(x)=Z_A+i_1x
\]
\[
 Z_{\mathrm{tang},B}(x)=Z_B+i_2(x-L)
\]

com
\[
 Z_B=Z_A+\frac{(i_1+i_2)L}{2}.
\]

Igualando:
\[
 Z_A+i_1x=Z_B+i_2(x-L)
\]

\[
\begin{aligned}
Z_A+i_1x
&=Z_A+\frac{(i_1+i_2)L}{2}+i_2x-i_2L \\
i_1x-i_2x
&=\frac{(i_1+i_2)L}{2}-i_2L \\
(i_1-i_2)x
&=\frac{(i_1+i_2-2i_2)L}{2} \\
(i_1-i_2)x
&=\frac{(i_1-i_2)L}{2} \\
x&=\frac{L}{2}.
\end{aligned}
\]

A cota no ponto de interseção é:
\[
 Z_I=Z_A+i_1\frac{L}{2}.
\]

Sabendo que a flecha vertical é
\[
 e=\frac{L}{8}(i_1-i_2),
\]

a cota da parábola no ponto $x=L/2$ é:
\[
 Z_{\mathrm{parab}}\left(\frac{L}{2}\right)=Z_A+i_1\frac{L}{2}-e.
\]

Assim, o ponto de interseção das tangentes ocorre em $x=L/2$, e a cota correspondente sobre a parábola é
\[
 Z=Z_A+i_1\frac{L}{2}-e.
\]
\end{corpo}

\section*{Parábola não simétrica em relação ao comprimento horizontal}

\begin{corpo}
Nos desenvolvimentos anteriores, o PIV foi considerado no ponto médio da curva, isto é, com $x_I=L/2$. Em uma concordância vertical não simétrica, essa condição não é mais verdadeira. Nesse caso, definem-se dois comprimentos distintos:
\[
 l_1=\text{distância do PCV ao PIV}
\]
\[
 l_2=\text{distância do PIV ao PTV}
\]
\[
 L=l_1+l_2
\]

Assim, tomando a origem no PCV, a abscissa do PIV passa a ser:
\[
 x_I=l_1
\]

Portanto, nas equações da curva simétrica, não se deve substituir simplesmente $L$ por $l_1+l_2$ e manter $L/2$. O comprimento total continua sendo $L=l_1+l_2$, mas a posição do PIV deve ser tratada separadamente.
\end{corpo}

\subsubsection*{Cotas dos pontos extremos a partir do PIV}

\begin{corpo}
Se a cota do PIV é $Z_I$, a cota do PCV, indicado por $A$, é obtida pela tangente inicial:
\[
 Z_A=Z_I-i_1l_1
\]

A cota do PTV, indicado por $B$, é obtida pela tangente final:
\[
 Z_B=Z_I+i_2l_2
\]

Dessa forma, para a parábola assimétrica:
\[
 Z_A=Z_I-i_1l_1
\]
\[
 Z_B=Z_I+i_2l_2
\]
\end{corpo}

\subsubsection*{Demonstração da flecha na parábola assimétrica}

\begin{corpo}
Na curva simétrica, a flecha vertical é avaliada em $x=L/2$. Na curva não simétrica, a flecha deve ser avaliada na abscissa do PIV, isto é, em $x=l_1$.

Seja $s$ a declividade da parábola no ponto da curva situado sob o PIV. Como a concordância é composta por dois trechos parabólicos, a declividade nesse ponto deve ser a mesma quando se chega pelo lado do PCV e quando se sai pelo lado do PTV. Essa declividade intermediária é:
\[
 s=\frac{i_1l_1+i_2l_2}{L}
\]

No trecho entre o PCV e o PIV, a variação de cota da parábola no ponto $x=l_1$ é:
\[
 Z_\Gamma=Z_A+i_1l_1+\frac{s-i_1}{2l_1}l_1^2
\]

ou seja,
\[
 Z_\Gamma=Z_A+i_1l_1+\frac{(s-i_1)l_1}{2}.
\]

Como $Z_A=Z_I-i_1l_1$, resulta:
\[
 Z_\Gamma=Z_I+\frac{(s-i_1)l_1}{2}.
\]

Substituindo $s$:
\[
 s-i_1=\frac{i_1l_1+i_2l_2}{L}-i_1
\]
\[
 s-i_1=\frac{i_1l_1+i_2l_2-i_1(l_1+l_2)}{L}
\]
\[
 s-i_1=\frac{(i_2-i_1)l_2}{L}=-\frac{(i_1-i_2)l_2}{L}
\]

Logo:
\[
 Z_\Gamma=Z_I-\frac{(i_1-i_2)l_1l_2}{2L}
\]

A flecha vertical $e$ é a diferença entre o PIV e o ponto correspondente da curva, $\Gamma$:
\[
 e=Z_I-Z_\Gamma=\frac{l_1l_2(i_1-i_2)}{2L}
\]
\end{corpo}

\subsubsection*{Relação com a forma simétrica}

\begin{corpo}
A expressão da flecha na parábola assimétrica é:
\[
 e=\frac{l_1l_2(i_1-i_2)}{2L}
\]

Quando a curva é simétrica, tem-se:
\[
 l_1=l_2=\frac{L}{2}.
\]

Substituindo na expressão anterior:
\[
\begin{aligned}
e
&=\frac{\left(\frac{L}{2}\right)\left(\frac{L}{2}\right)(i_1-i_2)}{2L} \\
&=\frac{L^2/4}{2L}(i_1-i_2) \\
&=\frac{L(i_1-i_2)}{8}.
\end{aligned}
\]

Portanto, a equação usual da flecha é um caso particular da equação mais geral:
\[
 e=\frac{L(i_1-i_2)}{8}\quad \text{quando}\quad l_1=l_2=\frac{L}{2}.
\]
\end{corpo}

\subsubsection*{Interpretação do sinal de $e$}

\begin{corpo}
Mantendo as rampas em forma decimal:
\begin{itemize}
    \item se $i_1>i_2$, a curva é convexa e $e$ é positivo, indicando que a parábola fica abaixo do PIV;
    \item se $i_1<i_2$, a curva é côncava e $e$ é negativo pela convenção algébrica;
    \item quando se deseja apenas a distância geométrica, utiliza-se $|e|$.
\end{itemize}
\end{corpo}

\subsubsection*{Equações da parábola composta por trechos}

\begin{corpo}
A parábola não simétrica deve ser escrita por trechos, pois os comprimentos entre PCV--PIV e PIV--PTV são distintos. Seja:
\[
 L=l_1+l_2 \quad \text{e} \quad s=\frac{i_1l_1+i_2l_2}{L}
\]

\paragraph*{Trecho do PCV ao ponto sob o PIV}

Para $0\leq x\leq l_1$, com origem no PCV:
\[
 Z(x)=Z_A+i_1x+\frac{s-i_1}{2l_1}x^2
\]

Como:
\[
 s-i_1=-\frac{(i_1-i_2)l_2}{L}
\]

então:
\[
 Z(x)=Z_A+i_1x-\frac{(i_1-i_2)l_2}{2Ll_1}x^2;\quad 0\leq x\leq l_1
\]

\paragraph*{Trecho do ponto sob o PIV ao PTV}

Para $l_1\leq x\leq L$, escrevendo a equação a partir do PTV:
\[
 Z(x)=Z_B+i_2(x-L)+\frac{i_2-s}{2l_2}(x-L)^2
\]

Como:
\[
 i_2-s=-\frac{(i_1-i_2)l_1}{L}
\]

então:
\[
 Z(x)=Z_B+i_2(x-L)-\frac{(i_1-i_2)l_1}{2Ll_2}(x-L)^2;\quad l_1\leq x\leq L
\]

Essas equações garantem a tangência às rampas nos extremos e a continuidade da cota e da declividade no ponto da curva situado sob o PIV.
\end{corpo}

\subsubsection*{Exemplo numérico}

\begin{corpo}
Considere uma curva vertical assimétrica com os seguintes dados:
\begin{itemize}
    \item $i_1=+3{,}0\%=0{,}030$;
    \item $i_2=-1{,}0\%=-0{,}010$;
    \item $l_1=80\,\text{m}$;
    \item $l_2=120\,\text{m}$;
    \item $Z_I=500{,}000\,\text{m}$.
\end{itemize}

O comprimento total da curva é:
\[
 L=l_1+l_2=80+120=200\,\text{m}
\]

A cota do PCV é:
\[
 Z_A=Z_I-i_1l_1=500{,}000-0{,}030\cdot 80=497{,}600\,\text{m}
\]

A cota do PTV é:
\[
 Z_B=Z_I+i_2l_2=500{,}000+(-0{,}010)\cdot 120=498{,}800\,\text{m}
\]

A declividade comum no ponto da curva situado sob o PIV é:
\[
 s=\frac{0{,}030\cdot 80+(-0{,}010)\cdot 120}{200}=0{,}006=0{,}6\%.
\]

A flecha vertical é:
\[
 e=\frac{80\cdot 120\,[0{,}030-(-0{,}010)]}{2\cdot 200}
\]
\[
 e=\frac{9600\cdot 0{,}040}{400}=0{,}960\,\text{m}
\]

Logo, a cota da curva sob o PIV é:
\[
 Z_\Gamma=Z_I-e=500{,}000-0{,}960=499{,}040\,\text{m}
\]
\end{corpo}

\subsubsection*{Equações do exemplo}

\begin{corpo}
Para o trecho entre o PCV e o ponto sob o PIV:
\[
 Z(x)=497{,}600+0{,}030x-\frac{(0{,}030-(-0{,}010))\cdot 120}{2\cdot 200\cdot 80}x^2
\]
\[
 Z(x)=497{,}600+0{,}030x-0{,}000150x^2
\]

Para $x=80\,\text{m}$:
\[
\begin{aligned}
Z(80)&=497{,}600+0{,}030\cdot 80-0{,}000150\cdot 80^2 \\
&=497{,}600+2{,}400-0{,}960 \\
&=499{,}040\,\text{m}.
\end{aligned}
\]

Para o trecho entre o ponto sob o PIV e o PTV:
\[
 Z(x)=498{,}800-0{,}010(x-200)-\frac{(0{,}030-(-0{,}010))\cdot 80}{2\cdot 200\cdot 120}(x-200)^2
\]
\[
 Z(x)=498{,}800-0{,}010(x-200)-0{,}000066667(x-200)^2
\]

Para $x=80\,\text{m}$:
\[
\begin{aligned}
Z(80)&=498{,}800-0{,}010(80-200)-0{,}000066667(80-200)^2 \\
&=498{,}800+1{,}200-0{,}960 \\
&=499{,}040\,\text{m}.
\end{aligned}
\]

As duas equações fornecem a mesma cota no ponto sob o PIV, confirmando a continuidade da concordância:
\[
 Z_\Gamma=499{,}040\,\text{m}
\]

Além disso, a declividade nesse ponto é a mesma nos dois trechos:
\[
 s=0{,}006=0{,}6\%.
\]
\end{corpo}

\end{document}
